\section{Related Work}
%==============================================================================
We organise prior art into themes.

\textbf{Transparent and glass-like segmentation.} Trans10K and TransLab
\cite{xie2020translab} established the benchmark and a boundary-aware baseline;
Trans2Seg \cite{xie2021trans2seg} introduced a transformer decoder. Glass and
mirror detection has explored boundary and reflection priors: enhanced boundary
learning \cite{he2021eblnet}, large-field glass detection \cite{mei2020gdnet},
mirror-specific cues \cite{yang2019mirrornet}, reflection-prior context
aggregation \cite{lin2021gsd}, polarization cues
\cite{mei2022polarization}, and video glass surface detection from motion
\cite{liu2024vgsd}. Our \brf{} module belongs to this boundary-prior
family; our contribution is not a new prior but a controlled measurement of its
(lack of) effect atop a strong backbone. Robotics-oriented transparent datasets
such as ClearGrasp \cite{sajjan2020cleargrasp}, ClearPose
\cite{chen2022clearpose}, and TransCG \cite{fang2022transcg} target depth/pose;
we use ClearPose for cross-dataset segmentation transfer.

\textbf{Segmentation architectures.} Encoder--decoder and dilated-convolution
designs---U-Net \cite{ronneberger2015unet}, DeepLabV3+
\cite{chen2018deeplabv3plus}, PSPNet \cite{zhao2017pspnet}, HRNet
\cite{wang2020hrnet}---and multi-scale fusion \cite{lin2017fpn} form classical
baselines. Vision transformers \cite{dosovitskiy2021vit} and hierarchical
variants \cite{liu2021swin} underpin transformer segmenters: SETR
\cite{zheng2021setr}, SegFormer \cite{xie2021segformer}, Mask2Former
\cite{cheng2022mask2former}, and UPerNet \cite{xiao2018upernet}.

\textbf{Self-supervised and foundation models.} DINO \cite{caron2021dino} and
\dinov{} \cite{oquab2023dinov2} yield features sensitive to surface boundaries;
masked \cite{he2022mae} and contrastive \cite{he2020moco} pretraining and
language supervision \cite{radford2021clip} are complementary. Promptable
foundation segmenters---SAM \cite{kirillov2023sam}, SAM~2
\cite{ravi2024sam2}, and HQ-SAM \cite{ke2023hqsam}---perform poorly zero-shot on
transparency, as we confirm.

\textbf{Optical flow and physics-guided learning.} \ofcv{} builds on RAFT
\cite{teed2020raft}; alternatives include PWC-Net \cite{sun2018pwcnet}, FlowNet
2.0 \cite{ilg2017flownet2}, and GMFlow \cite{xu2022gmflow}. Our cross-attention
follows the transformer formulation \cite{vaswani2017transformer} with
channel/spatial attention \cite{woo2018cbam,hu2018senet}. Optional graph
refinement uses SLIC superpixels \cite{achanta2012slic} and graph networks
\cite{kipf2017gcn,xu2019gin}. Training uses Dice supervision
\cite{milletari2016vnet}, AdamW \cite{loshchilov2019adamw}, and cosine schedules
\cite{loshchilov2017sgdr}. We evaluate robustness with common corruptions
\cite{hendrycks2019robustness}, interpretability with Grad-CAM
\cite{selvaraju2017gradcam}, and frame our study against physics-informed
learning \cite{karniadakis2021piml}, to which it is a cautionary data point.

%==============================================================================
